%% ==========================================================================
%% SECTION 6: RELATED WORK
%% ==========================================================================
\section{Related Work}\label{sec:related}

Our contribution is a formal impossibility result. Prior work documents
symptoms; we identify the cause. Prior work proposes treatments; we explain why
they fail to cure. The organizing principle is the distinction between
\emph{observation} (hallucination happens, training helps sometimes) and
\emph{explanation} (text-only observation plus bounded supervision makes
epistemic honesty unverifiable).

\fakepara{Hallucination taxonomies.} Extensive empirical work catalogs LLM
fabrication patterns: citation hallucination in academic
contexts~\cite{alkaissi2023artificial}, medical
misinformation~\cite{thirunavukarasu2023large}, and legal
fabrication~\cite{dahl2024large}. These findings instantiate our
\autoref{cond:hallucination}: in each domain, plausible fabrications are
cheap to generate and expensive to verify, placing supervisors in the hallucination
regime where gradient signal cannot distinguish truth from fabrication. The
taxonomies document the symptom; our theorems identify the structural cause.

\fakepara{Activation interpretation.} Recent work on Activation Oracles
demonstrates that LLMs can be trained to interpret the internal states of other models,
translating activation patterns into natural language explanations~\cite{karvonen2025activation}. This approach
addresses State Exteriority by making internal states available to an external system.
However, the authors note that the oracle "frequently makes incorrect guesses" and "is not
trained to express uncertainty." If the oracle's outputs are evaluated only as text by
bounded supervisors, the verification problem recurs at the meta-level. Full escape
requires that the oracle's claims be independently verifiable against the activation
data, which is a form of Provenance Binding applied to the interpretation process itself.


\fakepara{RLHF critiques.} Work on sycophancy~\cite{perez2022discovering},
reward hacking~\cite{gao2023scaling}, and the tension between helpfulness and
honesty~\cite{bai2022training} demonstrates that preference optimization can
degrade factual accuracy. Our \autoref{thm:learnability} formalizes why:
bounded supervisors cannot distinguish plausible fabrications from truthful
responses, so the learning signal is systematically corrupted.

\fakepara{Retrieval augmentation.} RAG systems~\cite{lewis2020retrieval} and
grounded generation approaches attempt to escape fabrication by conditioning on
external documents. Our analysis shows this is necessary but not sufficient:
State Exteriority addresses \autoref{thm:representational}, but without
Verification Independence, the learning problem
(\autoref{thm:learnability}) persists.

\fakepara{Uncertainty quantification.} Calibration methods, conformal
prediction~\cite{angelopoulos2021gentle}, and verbalized
uncertainty~\cite{kadavath2022language} attempt to surface model confidence. Our
results explain why these approaches are fundamentally limited under text-only
interfaces: the observation model discards the internal state that would
distinguish calibrated uncertainty from confident fabrication.

\fakepara{Semantic entropy.} Kuhn et
al.~\cite{kuhn2023semanticuncertaintylinguisticinvariances} propose sampling
multiple responses and clustering them by meaning to estimate uncertainty,
achieving strong results on question-answering benchmarks. Their method is the
most sophisticated text-channel uncertainty approach to date: by measuring
agreement \emph{between} outputs rather than features \emph{within} a single
output, it partially addresses the gaming problem. However, it remains a
text-channel method: the clustering operates on generated text, the uncertainty
estimate is derived from output diversity, and a bounded supervisor evaluating
semantic equivalence is subject to the same observation limits our theorems
describe. Our approach differs structurally: rather than sampling multiple
outputs and measuring their agreement (costly and still text-observable), we
extract signals from the computation that produced a single output---signals
the model cannot independently control. The two approaches are complementary:
semantic entropy measures output-level consistency; the tensor interface measures
computation-level confidence.

\fakepara{Impossibility results in LLMs.} Xu et
al.~\cite{xu2024hallucination} prove that hallucination is inevitable for LLMs
over the class of computable functions, which is a computational limit. Our result is
orthogonal: we prove that even if a system could be honest, a bounded
supervisor observing only text cannot verify that honesty, which is an observational
limit. Their impossibility concerns what LLMs can compute; ours concerns what
supervisors can verify. Both are structural, but they constrain different
dimensions of system design.

\fakepara{Impossibility results in distributed systems.} The
Fischer-Lynch-Paterson theorem~\cite{flp} proved that consensus is impossible
under asynchrony with one faulty process. This did not end distributed systems
research; it clarified the design space. Paxos~\cite{paxos}, Raft, and modern
consensus protocols work \emph{around} FLP by relaxing assumptions or
strengthening guarantees. Our result plays an analogous role: text-only
observation models with bounded supervision cannot guarantee epistemic honesty,
but systems that export structured epistemic state can escape the impossibility
class.

\fakepara{Epistemic evaluation frameworks.} Recent work proposes benchmarks
and frameworks for measuring epistemic properties of language
models~\cite{wang2024epistemic}. The Epistemic Alignment framework explicitly
notes that current interfaces do not let users specify uncertainty, sourcing,
or perspectives in structured ways, and that providers lack verification tools.
These critiques align with our analysis but stop at interface desiderata and
policy analysis. We provide the formal impossibility that explains \emph{why}
text-only evaluation cannot suffice, and the tensor interface that demonstrates
a concrete escape.

\fakepara{Epistemic agent architectures.} Work on structuring epistemic
integrity in reasoning systems proposes architectures with explicit
propositions, justification graphs, contradiction detection, and chain-of-reason
logging~\cite{chen2025epistemic}. These proposals are closest in spirit to our
escape conditions: they recognize that epistemic honesty requires structured
internal state. However, they are typically framed as \emph{replacements} for
stochastic language generation as logically grounded agents rather than wrappers
around existing LLMs. Our contribution differs in two ways: we derive the
tensor interface from a formal impossibility result (not as a standalone design
choice), and we demonstrate that the escape is achievable as a minimal
\emph{addition} to current architectures, not a wholesale replacement.

The gap we fill is formalization. Prior work documents that LLMs fabricate (the
symptom), critiques training methods (proposed treatments), and proposes
mitigations (partial remedies). We prove \emph{why} the problem is structural
(the diagnosis) and \emph{what} architectural properties are required to escape
it (the cure). The symptom/cause distinction is not merely rhetorical: it
determines whether improvements are incremental patches or genuine escapes from
the impossibility class.
